\documentclass[a4paper, 12pt]{article}

\usepackage[english, norsk]{babel}
\usepackage{csquotes}
\usepackage[backend=biber, style=ieee]{biblatex}
\addbibresource{ref.bib}
\usepackage[margin=25mm]{geometry}
\usepackage[table]{xcolor}
\usepackage{graphicx}
\usepackage{url}
\usepackage{hyperref}
\usepackage{caption}
\usepackage{float}
\usepackage{changepage,threeparttable}
\usepackage{fancyhdr}
\headheight=15pt

\begin{document}

\begin{titlepage}

    \begin{center}

        {\LARGE TDT4140 - Programvareutvikling\\[0.25em]Leveranse 8}
        \\[4em]

        {
        Antall ord utenom forside og referanser:
        \\1351
        }

        \vfill
        
        \includegraphics[width=8cm]{ntnu.pdf}
        
        
    \end{center}
    
    \thispagestyle{empty}
    \addtocounter{page}{-1}

\end{titlepage}
    

\tableofcontents
\thispagestyle{empty}
\addtocounter{page}{-1}
\newpage

\pagestyle{fancy}
\fancyhf{}
\fancyhead[L]{Tema: 4 - Parprogrammering}
\fancyhead[R]{Antall ord: \{ ord \}}
\fancyfoot[C]{\thepage}


% Hva er parprogrammering?
% Presentasjon av problemstilling
% Hvordan bruke parprogrammering i praksis?
% - Basere på erfaringer fra prosjektet

\section{Hva er parprogrammering?}

% presentere teori og definisjoner av parprogrammering
% presentere min definisjon av parprogrammering
% 

Parprogrammering er en metode innen ekstremprogrammering der to utviklere jobber sammen. Den ene utvikleren, <<sjåføren>>, skriver koden mens den andre, <<navigatøren>>, følger med, leser koden, og gir tilbakemeldinger underveis. Etter en bestemt tid kan utviklere bytte roller \cite{McDowell2006}. 
I en meta-analyse gjort av Hannay, Dybå, Arisholm og Sjøberg i 2009 kommer det frem at parprogrammering er et nyttig verktøy for å øke kvaliteten på komplekse oppgaver, og produktivitet på enklere oppgaver \cite{hannay2009}.
En annen studie, av McDowell et. al. i 2006 viser at parprogrammering har stor virkning som et pedagogisk verktøy, og at det er en god måte for studenter å lære programmering på \cite{McDowell2006}.
Parprogrammering ser også ut til å være nyttig for å ivareta og utvikle studenters interesse for programmering, og særlig blant kvinner, en stadig underrepresentert gruppe i IT-bransjen \cite{McDowell2006}\cite{kattan2018}.

Parprogrammering differensierer seg fra å bare <<kode-sammen>> ved at det er en strukturert måte å jobbe på. Den ene utvikleren er den som skriver koden, mens den andre utvikleren er den som følger med og gir tilbakemeldinger \cite{kattan2018}. 
Dette er en viktig forskjell fra å bare kode sammen, hvor det er vanskeligere å få tilbakemeldinger og gi tilbakemeldinger på koden.

\section{Problemstilling}

% Hvordan påvirker parprogrammering samarbeid og produktivitet i et utviklingsteam? En sammenlikning av teori og erfaringer fra et prosjekt.

I denne oppgaven har jeg valg å ta for meg problemstillingen <<Hvordan påvirker parprogrammering samarbeid og produktivitet i et utviklingsteam?>>.
Oppgaver omfavner problemstillingen i lys av teori og erfaringer fra et prosjekt.

\section{Hvordan bruke parprogrammering i praksis?}

\newpage
\printbibliography[heading=bibintoc]

\end{document}
